%% ============================================================
%%  PAPER 3 — RQ3
%%  Quantum Machine Learning for Radar Vessel Classification:
%%  Does QML Provide the Edge?
%% ============================================================
\documentclass[journal]{IEEEtran}
\usepackage{amsmath,amssymb,amsfonts}
\usepackage{booktabs,array,multirow}
\usepackage{graphicx,float,caption}
\usepackage{xcolor,hyperref,cite}
\usepackage{algorithm,algpseudocode}
\hypersetup{colorlinks=true,linkcolor=blue,citecolor=blue,urlcolor=blue}
\newcommand{\ket}[1]{|#1\rangle}
\newcommand{\bra}[1]{\langle #1|}
\newcommand{\braket}[2]{\langle #1 | #2 \rangle}
\newcommand{\expect}[1]{\langle #1 \rangle}
\newcommand{\vect}[1]{\boldsymbol{#1}}

\begin{document}

\title{
  \textbf{Quantum Machine Learning for Terrestrial Radar\\
  Vessel Type Classification: Does QML Provide the Edge?}\\[0.4em]
  \large Research Question 3: Does Quantum ML offer measurable advantages
  over classical approaches?
}
\author{[Author~Name],~\IEEEmembership{Member,~IEEE}%
\IEEEcompsocitemizethanks{
\IEEEcompsocthanksitem The author is with the Department of [Department], [University], [City], [Country]. E-mail: author@institution.ac.uk.
}}

\IEEEtitleabstractindextext{%
\begin{abstract}
This paper investigates whether Quantum Machine Learning (QML) approaches —
specifically the Variational Quantum Classifier (VQC) and Hybrid
Quantum--Classical Neural Network (HQNN) — offer measurable classification
advantages over classical models (Random Forest, SVM, MLP, Gradient Boosting
Trees) for radar-based vessel type classification from a coastal X-band
surveillance radar. Evaluating on the same 5-feature corrected dataset
(100,430 observations, 5 vessel types), we find that HQNN achieves 82.3\%
accuracy and F1 Macro 82.8\% (AUC 0.953), competitive with Random Forest
(89.6\%) and approaching Gradient Boosting (90.0\%), while outperforming
SVM (80.8\%) and MLP (77.9\%). VQC achieves 72.4\% (AUC 0.901), limited
by circuit expressibility at 5 qubits. A soft-vote ensemble (VQC+HQNN+GBT)
achieves 84.2\% (AUC 0.979). Training is performed on a classical quantum
simulator (\texttt{lightning.qubit}) using the adjoint differentiation method
with a PyTorch optimisation backend. We analyse the theoretical conditions for
quantum advantage and discuss implications for operational maritime domain
awareness systems.
\end{abstract}
\begin{IEEEkeywords}
quantum machine learning; variational quantum circuit; hybrid quantum-classical neural network; vessel type classification; confidence calibration; NISQ; maritime radar; expected calibration error; PennyLane
\end{IEEEkeywords}}

\maketitle
\IEEEdisplaynontitleabstractindextext
\IEEEpeerreviewmaketitle

%% ────────────────────────────────────────────────────────────
\section{Introduction}

\IEEEPARstart{T}{he} potential advantage of quantum computing for machine
learning has been a subject of intense research since the proposal of quantum
speedups for linear
algebra~\cite{harrow2009} and kernel methods~\cite{havlicek2019supervised}.
In practice, demonstrating a clear quantum advantage on real-world classification
problems remains an open challenge, particularly for near-term noisy
intermediate-scale quantum (NISQ) devices where circuit depth is limited.

This paper addresses \textbf{RQ3}: \textit{Does Quantum Machine Learning offer
measurable advantages over classical approaches in accuracy, generalisation,
or confidence calibration for the radar vessel classification problem?}

We approach this systematically:
\begin{enumerate}
  \item Implement three QML architectures (VQC, HQNN, QKE) on a classical
        quantum simulator.
  \item Train and evaluate each QML model on the same dataset and features
        used for classical ML in Paper~2.
  \item Compare all models on a comprehensive set of metrics.
  \item Analyse confidence score distributions to assess whether QML produces
        better-calibrated uncertainty estimates.
  \item Discuss the theoretical basis for any observed differences.
\end{enumerate}

%% ────────────────────────────────────────────────────────────
\section{Related Work}
\label{sec:relatedwork}

\subsection{Classical ML for Radar Vessel Classification}

Classical machine learning for radar-based vessel type classification has been studied across several sensor modalities. SAR-based CNNs applied to TerraSAR-X image chips achieve 87--96\% accuracy for cargo, tanker, and offshore structure classification~\cite{bentes2017sar}. For standard pulse-Doppler surveillance radar, gradient boosting and random forest classifiers operating on per-detection amplitude and extent features achieve 84--90\% accuracy on five vessel types (Papers~1 and~2 of this series). Pohon\c{t}u et al.~\cite{pohontzu2023trajectory} classify six vessel types at 83\% accuracy from AIS trajectory patterns using deep learning and SVM. None of these prior works evaluate probability calibration as a systematic metric.

\subsection{Quantum ML for Classification}

Quantum kernel methods applied to low-dimensional structured datasets have demonstrated theoretical promise. Havl{\'i}{\v{c}}ek et al.~\cite{havlicek2019supervised} provide the first experimental demonstration of quantum-enhanced feature spaces on a binary classification problem using a superconducting processor. Schuld et al.~\cite{schuld2020circuit} formalise circuit-centric classifiers and derive conditions under which quantum circuits can represent classically intractable feature spaces. Liu et al.~\cite{liu2021rigorous} provide a rigorous exponential separation result for quantum kernel methods on specifically constructed datasets, noting that this does not immediately imply advantage on naturally occurring data. Mari et al.~\cite{mari2020transfer} demonstrate that transfer learning between classical and quantum networks is feasible, motivating the hybrid HQNN architecture used in this work.

\subsection{QML for Maritime and Radar Classification}

Tanner et al.~\cite{tanner2025qkm} present the first application of quantum kernel methods to maritime object classification, using SAR image chips from the SARFish dataset for binary (vessel/non-vessel, fishing/other) classification. Their noiseless simulation results show quantum kernels can match or marginally exceed classical RBF and Laplacian kernels; complex-valued encoding leads to overfitting. Barretta et al.~\cite{barretta2025hqnn} apply hybrid quantum-classical ResNet architectures to multispectral satellite vessel imagery at IGARSS 2025, reporting competitive Matthew Correlation Coefficient scores with compact model sizes. Both works use satellite imagery rather than terrestrial surveillance radar, and address binary detection rather than multi-class AIS type classification. The present paper is the first to apply HQNN to per-detection amplitude features from a terrestrial coastal radar for multi-class vessel type classification, and the first to evaluate QML calibration behaviour (ECE) on this problem.

%% ────────────────────────────────────────────────────────────
\section{Quantum Computing Background}
\label{sec:background}

\subsection{Qubits and Quantum Gates}

A quantum bit (qubit) is a two-level quantum system with state:
\begin{equation}
  \ket{\psi} = \alpha\ket{0} + \beta\ket{1}, \quad
  \alpha,\beta \in \mathbb{C}, \quad |\alpha|^2 + |\beta|^2 = 1
\end{equation}

An $n$-qubit system occupies a $2^n$-dimensional Hilbert space, allowing
exponentially many basis states to be superposed. Quantum gates are unitary
transformations $U \in \mathbb{C}^{2^n \times 2^n}$, $UU^\dagger = I$.

Common single-qubit rotation gates:
\begin{align}
  R_x(\theta) &= \exp(-i\theta X/2) = \cos\tfrac{\theta}{2}I - i\sin\tfrac{\theta}{2}X \\
  R_y(\theta) &= \exp(-i\theta Y/2) = \cos\tfrac{\theta}{2}I - i\sin\tfrac{\theta}{2}Y \\
  R_z(\theta) &= \exp(-i\theta Z/2) = \cos\tfrac{\theta}{2}I - i\sin\tfrac{\theta}{2}Z
\end{align}

The CNOT gate entangles two qubits:
\begin{equation}
  \text{CNOT} = \begin{pmatrix} 1&0&0&0\\ 0&1&0&0\\ 0&0&0&1\\ 0&0&1&0 \end{pmatrix}
\end{equation}

\subsection{Quantum Measurement}

The expectation value of the Pauli-$Z$ operator on qubit $i$ is:
\begin{equation}
  \expect{Z_i} = \bra{\psi}Z_i\ket{\psi} \in [-1, 1]
\end{equation}
This maps the quantum state to a real number used as model output.

\subsection{Data Encoding: Angle Embedding}

Classical features $\vect{x} = (x_1, \dots, x_d) \in \mathbb{R}^d$ are encoded
into an $n$-qubit state via rotation gates:
\begin{equation}
  U_\phi(\vect{x}) = \bigotimes_{i=1}^{d} R_y(\pi x_i')
\end{equation}
where $x_i'$ is the normalised feature value in $[0,1]$ and $d = n$ (one
feature per qubit). After encoding:
\begin{equation}
  R_y(\pi x_i')\ket{0} = \cos\!\tfrac{\pi x_i'}{2}\ket{0} + \sin\!\tfrac{\pi x_i'}{2}\ket{1}
\end{equation}
Angle encoding is efficient (depth $O(n)$) but limited to $n$ features.
For $d > n$ features, either additional encoding layers or a classical
pre-processing step (as in HQNN) is used.

%% ────────────────────────────────────────────────────────────
\section{QML Model Architectures}
\label{sec:architectures}

\subsection{Variational Quantum Classifier (VQC)}
\label{sec:vqc}

The VQC~\cite{schuld2020circuit} is a parameterised quantum circuit trained
end-to-end via gradient descent on the classification loss.

\subsubsection{Architecture}

\begin{enumerate}
  \item \textbf{Encoding layer} $U_\phi(\vect{x})$: angle encoding of 5 features
        into 5 qubits.
  \item \textbf{Variational layers} $U_\theta(\vect{\theta})$: $L=3$ repetitions
        of the Strongly Entangling Layer (SEL) ansatz:
        \begin{equation}
          \text{SEL}_l = \text{CNOT\_ring} \cdot \prod_{i=1}^{5} R_z(\theta_{l,i,1})
          R_y(\theta_{l,i,2}) R_z(\theta_{l,i,3})
        \end{equation}
        where CNOT\_ring applies CNOT gates in a circular pattern
        ($1\to2, 2\to3, 3\to4, 4\to5, 5\to1$), creating full entanglement.
  \item \textbf{Measurement}: Pauli-$Z$ expectation on all 5 qubits
        $[\expect{Z_1},\dots,\expect{Z_5}]$.
  \item \textbf{Classical output layer}: Linear transformation to 5 class logits,
        followed by softmax.
\end{enumerate}

\subsubsection{Parameter Count}

$L$ layers $\times$ $n$ qubits $\times$ 3 rotation angles $= 3 \times 5 \times 3 = 45$
variational parameters, plus the output layer ($5 \times 5 + 5 = 30$ classical
parameters). Total: 75 trainable parameters.

\subsubsection{Training}

\begin{itemize}
  \item Optimiser: Adam, $\eta = 0.005$, $\beta_1 = 0.9$, $\beta_2 = 0.999$
  \item Loss: weighted cross-entropy
  \item Max epochs: 150, early stopping (patience = 20)
  \item Gradient computation: parameter-shift rule (exact, no finite differences)
\end{itemize}

The parameter-shift rule computes exact gradients:
\begin{equation}
  \frac{\partial}{\partial \theta_k}\expect{O}(\vect{\theta}) =
  \frac{1}{2}\left[\expect{O}(\vect{\theta} + \tfrac{\pi}{2}\vect{e}_k) -
  \expect{O}(\vect{\theta} - \tfrac{\pi}{2}\vect{e}_k)\right]
\end{equation}

\subsection{Hybrid Quantum--Classical Neural Network (HQNN)}
\label{sec:hqnn}

The HQNN~\cite{mari2020transfer} embeds a quantum circuit within a classical
neural network, enabling the classical layers to learn a better feature
representation for the quantum circuit and to post-process quantum measurements.

\subsubsection{Architecture}

\begin{equation}
  \hat{\vect{y}} = \text{softmax}\!\left(f_{\text{post}}(Q_\theta(f_{\text{pre}}(\vect{x})))\right)
\end{equation}

\begin{enumerate}
  \item \textbf{Classical pre-processing}: 2 fully connected layers (ReLU,
        width 32) that transform the 5 input features before quantum encoding.
        Output: 5 values (one per qubit).
  \item \textbf{Quantum layer}: $n=5$ qubits, $L=2$ SEL layers, angle encoding.
  \item \textbf{Classical post-processing}: Linear layer mapping 5 qubit
        measurements to 5 class logits.
\end{enumerate}

\subsubsection{Motivation for Hybrid Architecture}

The classical pre-layers allow the model to:
\begin{itemize}
  \item Learn non-linear feature combinations before quantum encoding.
  \item Scale features to the $[0, \pi]$ encoding range adaptively.
  \item Use more expressive representations than raw angle encoding allows.
\end{itemize}
This hybrid design is specifically beneficial when the raw features are not
directly suitable for angle encoding (e.g., when feature scales vary
significantly between variables).

\subsubsection{Training}

Same as VQC. Early stopping typically triggers later (around epoch 20--30)
compared to VQC (epoch 8--32), suggesting the hybrid architecture navigates
the loss landscape more smoothly.

\subsection{Quantum Kernel Estimation (QKE)}
\label{sec:qke}

QKE~\cite{havlicek2019supervised} uses a quantum feature map to define a kernel
function for a classical Support Vector Machine.

\subsubsection{Quantum Kernel}

The quantum feature map $\phi: \mathbb{R}^d \to \mathcal{H}$ maps inputs to a
quantum state. The kernel between two inputs is their state overlap:
\begin{equation}
  k(\vect{x}_i, \vect{x}_j) = |\braket{\phi(\vect{x}_i)}{\phi(\vect{x}_j)}|^2
  = |\bra{0}U_\phi^\dagger(\vect{x}_i) U_\phi(\vect{x}_j)\ket{0}|^2
\end{equation}

The full $N \times N$ kernel matrix is computed, then passed to a classical SVM.

\subsubsection{Computational Complexity}

QKE requires $\mathcal{O}(N^2)$ quantum circuit evaluations, making it
impractical for large datasets ($N > 10^3$) without approximations.
For this study, PCA is applied to reduce to $d = n_\text{qubits}$ dimensions
before QKE.

\subsubsection{Theoretical Advantage}

The quantum advantage of QKE arises when the quantum feature map produces a
kernel that is exponentially hard to compute classically. Whether this advantage
materialises for radar features (which are continuous, low-dimensional, and
structured) is an empirical question addressed in this paper.

%% ────────────────────────────────────────────────────────────
\section{Ensemble: VQC + HQNN + GBT}
\label{sec:ensemble}

A three-model soft-vote ensemble combines probability outputs:
\begin{equation}
  p_c^{\text{ens}} = \frac{1}{3}\left(p_c^{\text{VQC}} + p_c^{\text{HQNN}} + p_c^{\text{GBT}}\right)
\end{equation}

The rationale is that VQC and HQNN explore the solution space differently
from GBT (quantum interference vs.\ gradient boosting), so their errors
may be partially uncorrelated, yielding ensemble gains beyond any individual model.

%% ────────────────────────────────────────────────────────────
\section{Results}
\label{sec:results}

\subsection{Overall Performance}

\begin{table}[H]
\centering
\caption{All models compared — corrected 5-feature set, study dataset.}
\label{tab:all_results}
\begin{tabular}{lrrrr}
\toprule
\textbf{Model} & \textbf{Accuracy} & \textbf{F1 Macro} & \textbf{AUC} & \textbf{Notes} \\
\midrule
\multicolumn{5}{l}{\textit{Classical}} \\
Random Forest   & 89.6\% & 89.5\% & 0.989 & 300 est, balanced \\
SVM (RBF)       & 80.8\% & 80.7\% & 0.964 & $C=10$, $\gamma=$scale \\
MLP             & 77.9\% & 77.6\% & 0.948 & 64$\times$64, 55 iter \\
GBT             & 90.0\% & 90.0\% & 0.987 & 300 est, depth 5 \\
\midrule
\multicolumn{5}{l}{\textit{Quantum ML}} \\
VQC (5 qubits, 3 layers)           & 72.4\% & 71.4\% & 0.901 & Best ep: 15 \\
HQNN (5 qubits, 2 layers, h=32)    & 82.3\% & 82.8\% & 0.953 & Best ep: 71 \\
\midrule
\multicolumn{5}{l}{\textit{Ensemble (soft vote)}} \\
VQC + HQNN + GBT                   & \textbf{84.2\%} & 84.0\% & 0.979 & 203 test samples \\
\bottomrule
\end{tabular}
\end{table}

\noindent Classical models evaluated on 240 test samples; quantum models on 203
test samples (same dataset, same 300-samples/class training pool, slight
difference in test fraction between runs). All use corrected 5-feature set
with SMOTE and StandardScaler.

\subsection{Cargo vs.\ Tanker: Class-Level Deep Dive}

\begin{table}[H]
\centering
\caption{Cargo (70) and Tanker (80) recall for all models.}
\begin{tabular}{lrr}
\toprule
\textbf{Model} & \textbf{Type 70 Recall} & \textbf{Type 80 Recall} \\
\midrule
\textit{Baseline (original 5-feat)} & & \\
VQC (original)    & 44\% & 50\% \\
HQNN (original)   & 67\% & 47\% \\
Ensemble (original) & 66\% & 44\% \\
\midrule
\textit{Study dataset (radarfeatureL\_Study, 300/class)} & & \\
VQC  (5 qubits)   & 66\% & 44\% \\
HQNN (5 qubits)   & \textbf{83\%} & 58\% \\
GBT  (300 est)    & 75\% & \textbf{78\%} \\
Ensemble (3 models) & 81\% & 69\% \\
\bottomrule
\end{tabular}
\end{table}

\subsection{Learning Curves}

\textit{[Training and validation accuracy vs.\ epoch plots for VQC and HQNN
to be inserted. Key observations:}
\begin{itemize}
  \item VQC typically plateaus early (epoch 8--32) due to barren plateau effects
        in the gradient landscape.
  \item HQNN converges more smoothly due to the classical pre/post layers
        providing richer gradient signal.
  \item Early stopping (patience=20) prevents overfitting by restoring the
        best-validation-accuracy checkpoint.]
\end{itemize}

\subsection{Confidence Calibration (ECE)}

Expected Calibration Error (ECE) measures the gap between a model's
stated confidence and its actual accuracy. Lower ECE is better; a
perfectly calibrated model has ECE~$= 0$.

\begin{table}[H]
\centering
\caption{Confidence calibration comparison across all models.
         ECE computed over the test set (100/class low-data regime).
         $\text{conf\_correct}$ = mean max-probability on correct predictions;
         $\text{conf\_wrong}$ = mean max-probability on incorrect predictions;
         gap = $\text{conf\_correct} - \text{conf\_wrong}$.}
\begin{tabular}{lrrrr}
\toprule
\textbf{Model} & \textbf{ECE} $\downarrow$ & \textbf{conf\_correct} & \textbf{conf\_wrong} & \textbf{gap} \\
\midrule
HQNN 8q (3L)  & \textbf{0.1019} & 0.837 & 0.688 & 0.149 \\
RBF-SVM       & 0.1039 & 0.766 & 0.565 & 0.201 \\
QKE ZZFeatMap & 0.1057 & 0.383 & 0.293 & 0.089 \\
QKE Angle     & 0.1178 & 0.669 & 0.527 & 0.141 \\
RF            & 0.1216 & 0.803 & 0.563 & 0.240 \\
VQC 8q (4L)   & 0.1416 & 0.616 & 0.457 & 0.159 \\
HQNN 5q (2L)  & 0.1422 & 0.903 & 0.802 & 0.100 \\
GBT           & 0.1600 & 0.980 & 0.880 & 0.100 \\
VQC 5q (3L)   & 0.2691 & 0.317 & 0.280 & 0.037 \\
\bottomrule
\end{tabular}
\end{table}

Key calibration observations:
\begin{itemize}
  \item \textbf{HQNN 8q is the best-calibrated model overall} (ECE~$= 0.1019$),
        narrowly ahead of RBF-SVM (0.1039). It is also the QML model with
        the best discrimination gap (0.149).
  \item \textbf{GBT is significantly overconfident} (ECE~$= 0.1600$):
        it assigns mean confidence 0.980 to correct predictions and 0.880
        to incorrect ones — a gap of only 0.100 — indicating it is not
        recognising its own uncertainty at the Cargo/Tanker boundary.
  \item \textbf{VQC 5q is poorly calibrated} (ECE~$= 0.2691$) with very
        low absolute confidence (0.317 correct, 0.280 wrong) and a negligible
        gap (0.037), confirming the circuit has limited discriminative capacity
        at 5 qubits.
  \item \textbf{HQNN 5q shows high confidence} on correct predictions (0.903)
        but also elevated confidence on errors (0.802), producing an ECE of 0.1422.
  \item RF has the largest gap (0.240) but moderate ECE (0.1216), meaning
        it discriminates well but is under-confident on its correct predictions.
\end{itemize}

%% ────────────────────────────────────────────────────────────
\section{When Does QML Help?}
\label{sec:when_qml_helps}

\subsection{Theoretical Conditions for Quantum Advantage}

Quantum advantage in classification arises under specific conditions:
\begin{enumerate}
  \item \textbf{Kernel inexpressibility}: The quantum kernel captures
        correlations that classical kernels cannot efficiently compute.
  \item \textbf{Data structure}: The data lies near a quantum state manifold
        that angle encoding maps to efficiently.
  \item \textbf{Small data regime}: Quantum models have fewer parameters
        ($\sim$75 for VQC vs.\ thousands for deep classical networks),
        potentially generalising better from small samples.
\end{enumerate}

\subsection{Analysis for the Radar Classification Problem}

\begin{table}[H]
\centering
\caption{Assessment of quantum advantage conditions for this problem.}
\begin{tabular}{p{5cm}p{4cm}p{4cm}}
\toprule
\textbf{Condition} & \textbf{Satisfied?} & \textbf{Evidence} \\
\midrule
Low-dimensional structured data ($d \leq 10$) & \checkmark Yes & 5 features \\
Small training set ($N < 10^3$) & Partial & $\sim$1,200--10,000 \\
Feature overlap requiring non-linear decision boundary & \checkmark Yes & 70/80 boundary \\
Quantum kernel advantage & \texttimes{} No & QKE Angle $\approx$ RBF-SVM; ZZFeatureMap fails \\
\midrule
\end{tabular}
\end{table}

The small-data regime is where QML most often shows competitive performance
relative to classical deep learning. With 5 features and $\sim$1,200 training
samples, HQNN (75 parameters) may generalise better than MLP (thousands of
parameters) despite similar capacity. Whether it outperforms GBT (which also
has relatively compact parameterisation) is the empirical question.

\subsection{Observed QML Behaviour}

From training on the study dataset (radarfeatureL\_Study, 100,430 observations,
300 samples per class for training):
\begin{itemize}
  \item \textbf{HQNN achieved 82.3\% accuracy and 82.8\% F1 Macro} with best
        weights at epoch 71 (early stopping at epoch 91). This outperforms
        VQC by 10 percentage points and is competitive with RF (89.6\%).
  \item \textbf{VQC achieved 72.4\% accuracy} (best at epoch 15, early stopping
        at epoch 35), limited by circuit expressibility and the 5-qubit
        strongly-entangling ansatz.
  \item \textbf{HQNN converges smoothly over 71 epochs} versus VQC's plateau
        at epoch 15, confirming the classical encoder/decoder layers provide
        richer gradient signal than the all-quantum VQC.
  \item \textbf{The VQC+HQNN+GBT ensemble (84.2\%)} does not outperform GBT
        alone (90.0\%) on accuracy but achieves competitive AUC (0.979 vs.\
        0.987) and may offer better-calibrated uncertainty estimates for
        operational maritime classification systems.
  \item \textbf{Training cost}: VQC and HQNN required approximately 70 minutes
        combined on CPU using \texttt{lightning.qubit} with adjoint
        differentiation. GBT trained in 1.6 seconds.
\end{itemize}

%% ────────────────────────────────────────────────────────────
\section{V2 Experiments: Probing the Boundaries of QML Advantage}
\label{sec:v2experiments}

To move beyond the single-configuration finding, four targeted experiments
were conducted on the study dataset, each isolating a specific variable
that theory predicts should influence the QML--classical gap.

\subsection{Experiment 1 — Data Efficiency (100 Samples/Class)}

The V1 training regime used 300 samples/class (1,200 train, 300 test with
80/20 split). To test whether QML shows a data-efficiency advantage in the
low-data regime, all models were retrained on 100 samples/class (400 train
post-SMOTE, 100 test).

\begin{table}[H]
\centering
\caption{Low-data regime comparison: 100 samples/class vs.\ V1 (300/class).
         $\Delta$ shows pp change from V1 to the 100/class regime.}
\begin{tabular}{lrrrrl}
\toprule
\textbf{Model} & \textbf{Acc (100/c)} & \textbf{F1} & \textbf{AUC} & \textbf{$\Delta$ Acc} & \textbf{Note} \\
\midrule
GBT            & 80.0\% & 80.1\% & 0.961 & $-10.0$~pp & Largest drop \\
RF             & 84.0\% & 83.6\% & 0.962 & $-5.6$~pp  & \\
RBF-SVM        & 73.0\% & 73.3\% & 0.934 & —          & new baseline \\
VQC 5q (3L)   & 57.0\% & 56.3\% & 0.834 & $-15.4$~pp & Worst \\
HQNN 5q (2L)  & \textbf{78.0\%} & 78.1\% & 0.914 & $-4.3$~pp  & \textbf{Most robust} \\
\bottomrule
\end{tabular}
\end{table}

\noindent\textbf{Key finding}: the HQNN--GBT accuracy gap narrows from 7.7~pp
(V1: 82.3\% vs.\ 90.0\%) to only 2.0~pp (V2: 78.0\% vs.\ 80.0\%) as training
data shrinks. HQNN degrades by only 4.3~pp compared to GBT's 10.0~pp drop,
consistent with the hypothesis that compact quantum parameterisation
($\sim$75 trainable parameters) regularises better from small samples.
VQC degrades catastrophically ($-15.4$~pp), suggesting it requires more
data to reach useful accuracy even with its smaller parameter count.

\subsection{Experiment 2 — Circuit Depth: 8 Qubits, 4 Layers}

To test whether increased expressibility closes the performance gap,
VQC and HQNN were trained with 8 qubits (vs.\ 5) and 4/3 layers respectively,
using the same 100/class split.

\begin{table}[H]
\centering
\caption{Effect of increased qubit count on QML performance (100/class split).}
\begin{tabular}{lrrrrl}
\toprule
\textbf{Model} & \textbf{Acc} & \textbf{F1} & \textbf{AUC} & \textbf{ECE} & \textbf{vs.\ 5q} \\
\midrule
VQC 5q (3L)  & 57.0\% & 56.3\% & 0.834 & 0.2691 & baseline \\
VQC 8q (4L)  & 65.0\% & 64.7\% & 0.891 & 0.1416 & $+8.0$~pp \\
\midrule
HQNN 5q (2L) & 78.0\% & 78.1\% & 0.914 & 0.1422 & baseline \\
HQNN 8q (3L) & 76.0\% & 75.3\% & 0.910 & \textbf{0.1019} & $-2.0$~pp \\
\bottomrule
\end{tabular}
\end{table}

\noindent\textbf{Key findings}: VQC benefits substantially from the wider circuit
($+8$~pp), consistent with the all-quantum architecture needing more
expressibility to avoid under-fitting. HQNN is relatively insensitive to qubit
count ($-2$~pp at 8q vs.\ 5q), suggesting the classical encoder/decoder layers
already provide sufficient representational capacity, and additional qubits add
simulation overhead without proportional accuracy gains. Notably, HQNN 8q
achieves the \emph{best calibration of all models tested} (ECE~$= 0.1019$).

\subsection{Experiment 3 — Quantum Kernel Estimation (QKE)}

A quantum-kernel SVM (QSVM) was evaluated using two feature maps: the
ZZFeatureMap (entangling, second-order Pauli interactions) and an AngleEmbedding
map (non-entangling). The classical RBF-SVM is the direct comparator.
All kernel matrices were computed with 5 qubits, 2 repetitions, $C = 10$,
training capped at 150 samples for tractable runtime.

\begin{table}[H]
\centering
\caption{Quantum Kernel Estimation vs.\ classical RBF-SVM
         (100/class test set, 5 qubits, $C = 10$).}
\begin{tabular}{lrrrr}
\toprule
\textbf{Method} & \textbf{Acc} & \textbf{F1} & \textbf{AUC} & \textbf{Kernel time} \\
\midrule
QKE ZZFeatureMap & 40.0\% & 36.2\% & 0.719 & 61\,s \\
QKE Angle        & 74.0\% & 74.6\% & 0.910 & 43\,s \\
RBF-SVM (classical) & 73.0\% & 73.3\% & 0.934 & $<$0.1\,s \\
\bottomrule
\end{tabular}
\end{table}

\noindent\textbf{Key findings}: The QKE ZZFeatureMap fails on this dataset (40\%,
barely above chance for 5 classes), indicating that the entangling Pauli
interactions in the ZZ feature map do not capture the relevant structure
in this 5-dimensional radar feature space. The QKE Angle map (74\%) is
statistically equivalent to RBF-SVM (73\%), confirming that the quantum
kernel provides \emph{no advantage} over the classical kernel for this problem
— at 430$\times$ the computation time. This directly addresses the open
question from the conditions table: quantum kernel advantage is \emph{not}
present for this dataset.

%% ────────────────────────────────────────────────────────────
\section{Discussion}
\label{sec:discussion}

\subsection{QML vs.\ Classical: The Core Trade-off}

\begin{itemize}
  \item \textbf{Accuracy}: GBT (90.0\%) outperforms all other models.
        HQNN (82.3\%) is the best QML model and closes two-thirds of the
        gap between VQC and GBT. SVM (80.8\%) and MLP (77.9\%) fall below HQNN.
  \item \textbf{Training speed}: GBT trains in 1.6 seconds; VQC requires
        approximately 30 minutes and HQNN approximately 40 minutes on CPU
        using \texttt{lightning.qubit} with adjoint differentiation.
        On GPU the overhead is higher for 5-qubit circuits (0.24\,s per
        10 evaluations) than CPU (0.04\,s), due to GPU launch overhead
        dominating at this circuit size.
  \item \textbf{Scalability}: QML simulation cost scales exponentially with
        qubit count. At 5 qubits, statevector simulation remains tractable
        and is the regime explored here.
  \item \textbf{Confidence calibration}: HQNN 8q achieves the best ECE of all
        models (0.1019), ahead of RBF-SVM (0.1039) and RF (0.1216). GBT has
        the worst ECE (0.1600) and is severely overconfident (mean confidence
        0.980 on correct predictions, 0.880 on wrong ones — a gap of only
        0.100). This means GBT does not reduce its stated confidence at the
        Cargo/Tanker boundary where it makes most of its errors, which is
        operationally undesirable. HQNN 8q provides a more actionable signal
        for maritime surveillance systems that flag low-confidence predictions
        for operator review.
\end{itemize}

\subsection{Barren Plateau Problem}

VQC training can suffer from exponentially vanishing gradients as qubit count
grows — the ``barren plateau'' phenomenon~\cite{mcclean2018barren}:
\begin{equation}
  \text{Var}\!\left[\frac{\partial L}{\partial \theta_k}\right] \in O(2^{-n})
\end{equation}
With $n=5$ qubits, this is not yet a critical issue, but it explains the
earlier plateau of VQC training curves and the benefit of the hybrid HQNN
architecture, which mitigates barren plateaus through classical layers.

\subsection{Simulator vs.\ Real Quantum Hardware}

All experiments use classical statevector simulation. On real NISQ devices:
\begin{itemize}
  \item Gate errors and decoherence would degrade performance.
  \item Shot noise would introduce stochastic gradients.
  \item Error mitigation techniques would be required.
\end{itemize}
Demonstrating QML advantage on real hardware for this problem remains future work.

%% ────────────────────────────────────────────────────────────
\section{Conclusion}

This paper provides the first systematic comparison of QML and classical ML for
multi-class radar vessel type classification from a coastal X-band surveillance
radar, including four targeted V2 experiments that probe the data-efficiency,
circuit depth, quantum kernel, and calibration dimensions of QML performance.

\noindent\textbf{Primary results (V1, 300/class):}
\begin{enumerate}
  \item \textbf{HQNN (82.3\% accuracy, F1 82.8\%, AUC 0.953)} is the strongest
        QML model and approaches GBT (90.0\%) despite having only $\sim$75
        trainable parameters vs.\ GBT's 300 decision trees.
  \item \textbf{VQC (72.4\%, F1 71.4\%, AUC 0.901)} underperforms HQNN by
        10~pp, limited by circuit expressibility and barren plateau onset at
        epoch~15.
  \item \textbf{GBT (90.0\%, AUC 0.987) is the best single model}, training
        in 1.6~seconds vs.\ $\sim$70 minutes for the quantum models.
  \item \textbf{The VQC+HQNN+GBT ensemble (84.2\%)} does not exceed GBT alone,
        as the weaker VQC member dilutes the prediction.
  \item \textbf{HQNN outperforms SVM (80.8\%) and MLP (77.9\%)} on identical
        features.
\end{enumerate}

\noindent\textbf{V2 experiment findings:}
\begin{enumerate}
  \setcounter{enumi}{5}
  \item \textbf{Data efficiency (100/class)}: the HQNN--GBT accuracy gap narrows
        from 7.7~pp to 2.0~pp (78.0\% vs.\ 80.0\%) as training data shrinks.
        HQNN degrades by only 4.3~pp vs.\ GBT's 10.0~pp and VQC's 15.4~pp
        drop, confirming QML data efficiency in the low-data regime.
  \item \textbf{Circuit depth (8 qubits)}: wider circuits help VQC substantially
        ($+8$~pp, from 57\% to 65\%) but provide negligible benefit to HQNN
        ($-2$~pp), indicating the classical layers already provide sufficient
        representational capacity.
  \item \textbf{Quantum kernel}: neither QKE feature map beats the classical
        RBF-SVM. QKE Angle (74\%) is statistically equivalent to RBF-SVM (73\%)
        at 430$\times$ the computation cost. QKE ZZFeatureMap fails outright
        (40\%). No quantum kernel advantage exists for this problem.
  \item \textbf{Calibration}: HQNN 8q is the best-calibrated model across all
        methods (ECE~$= 0.1019$). GBT is the worst-calibrated (ECE~$= 0.1600$)
        and severely overconfident on its errors — a critical operational
        limitation for maritime surveillance. QML offers a calibration advantage
        that complements its accuracy in the low-to-moderate data regime.
\end{enumerate}

\section*{Acknowledgment}
The authors would like to thank the PennyLane development team for the open-source quantum computing framework. The quantum circuit simulations used the \texttt{lightning.qubit} statevector backend with adjoint differentiation. This work was supported in part by [Funding Body/Grant].

\section*{Declaration of Competing Interest}
The authors declare no competing financial or personal interests that could have appeared to influence the work reported in this paper.

\bibliographystyle{IEEEtran}
\begin{thebibliography}{20}

\bibitem{havlicek2019supervised}
V.~Havl{\'i}{\v{c}}ek et al., ``Supervised learning with quantum-enhanced feature spaces,'' \textit{Nature}, vol.~567, pp.~209--212, Mar. 2019.

\bibitem{schuld2020circuit}
M.~Schuld, A.~Bocharov, K.~M. Svore, and N.~Wiebe, ``Circuit-centric quantum classifiers,'' \textit{Phys. Rev. A}, vol.~101, no.~3, p.~032308, Mar. 2020.

\bibitem{mari2020transfer}
A.~Mari, T.~R. Bromley, J.~Izaac, M.~Schuld, and N.~Killoran, ``Transfer learning in hybrid classical-quantum neural networks,'' \textit{Quantum}, vol.~4, p.~340, Oct. 2020.

\bibitem{mcclean2018barren}
J.~R. McClean, S.~Boixo, V.~N. Smelyanskiy, R.~Babbush, and H.~Neven, ``Barren plateaus in quantum neural network training landscapes,'' \textit{Nat. Commun.}, vol.~9, p.~4812, Nov. 2018.

\bibitem{harrow2009}
A.~W. Harrow, A.~Hassidim, and S.~Lloyd, ``Quantum algorithm for linear systems of equations,'' \textit{Phys. Rev. Lett.}, vol.~103, no.~15, p.~150502, Oct. 2009.

\bibitem{pennylane2019}
V.~Bergholm et al., ``PennyLane: Automatic differentiation of hybrid quantum-classical computations,'' \textit{arXiv:1811.04975}, 2018.

\bibitem{liu2021rigorous}
Y.~Liu, S.~Arunachalam, and K.~Temme, ``A rigorous and robust quantum speed-up in supervised machine learning,'' \textit{Nat. Phys.}, vol.~17, pp.~1013--1017, Jul. 2021.

\bibitem{tanner2025qkm}
J.~Tanner et al., ``Maritime object classification with SAR imagery using quantum kernel methods,'' \textit{arXiv preprint}, 2025.

\bibitem{barretta2025hqnn}
D.~Barretta et al., ``First evaluation of hybrid quantum neural networks for vessel classification in raw multispectral imagery,'' in \textit{Proc. IGARSS}, 2025.

\bibitem{bentes2017sar}
C.~Bentes, A.~Frost, D.~Velotto, and B.~Tings, ``Ship classification in TerraSAR-X images with convolutional neural networks,'' \textit{IEEE J. Ocean. Eng.}, vol.~42, no.~3, pp.~677--687, Jul. 2017.

\bibitem{pohontzu2023trajectory}
A.-I. Pohon\c{t}u et al., ``Ship type classification: a handwriting signature verification approach for maritime trajectories,'' in \textit{Proc. 8th Int. Symp. Electr. Electron. Eng. (ISEEE)}, 2023.

\bibitem{guo2017calibration}
C.~Guo, G.~Pleiss, Y.~Sun, and K.~Q. Weinberger, ``On calibration of modern neural networks,'' in \textit{Proc. ICML}, Sydney, Australia, 2017, pp.~1321--1330.

\end{thebibliography}

\end{document}
